\section{METODOLOGIA}

O Capítulo 2 mostrou que a busca por um animal perdido depende de informações espalhadas por diferentes canais e de uma colaboração que se concentra na vizinhança. Para confirmar e detalhar essas necessidades no contexto do WeFIND, era preciso ouvir pessoas que já passaram por essa situação e examinar como as plataformas existentes tratam o problema. Por isso, o trabalho combinou uma pesquisa de campo com participantes da comunidade, uma análise comparativa de sistemas similares e um processo de desenvolvimento iterativo. Os procedimentos e resultados dessas etapas são apresentados a seguir; a avaliação de usabilidade, por envolver o aplicativo já implementado, é tratada no Capítulo 8.

\subsection{Classificação da Pesquisa}

Para estruturar as decisões tomadas ao longo do estudo, a pesquisa seguiu a classificação metodológica fundamentada em Gil (2022) e Marconi e Lakatos (2022), considerando sua natureza, objetivos, abordagem do problema e procedimentos técnicos empregados.

Quanto à sua natureza, a pesquisa classifica-se como aplicada. Como explica Gil (2022), a pesquisa aplicada busca produzir conhecimentos voltados à solução de problemas práticos. Neste trabalho, os conhecimentos levantados orientaram o projeto, a especificação e o desenvolvimento de uma aplicação móvel de apoio à divulgação e à busca por animais perdidos e encontrados.

Quanto aos seus objetivos, a pesquisa é exploratório-descritiva. O caráter exploratório está relacionado à investigação das necessidades e dificuldades associadas à busca por animais, enquanto o caráter descritivo se manifesta no registro das respostas dos participantes e das funcionalidades observadas nas plataformas analisadas (Gil, 2022; Marconi; Lakatos, 2022).

Quanto à abordagem, o estudo combina procedimentos quantitativos e qualitativos. Os dados quantitativos correspondem às respostas fechadas dos questionários, incluindo as dez afirmações baseadas na \textit{System Usability Scale} (SUS) e a conclusão autorrelatada das tarefas na avaliação do aplicativo. Para examinar as respostas abertas, realizou-se uma leitura das respostas e agruparam-se ideias recorrentes em temas identificados durante essa leitura. Essa organização teve caráter descritivo e apoiou a identificação de dificuldades, aspectos positivos e sugestões dos participantes (Gil, 2022; Sauro; Lewis, 2016).

Quanto aos procedimentos técnicos, o estudo reuniu pesquisa bibliográfica, pesquisa de campo e análise comparativa de sistemas similares (Gil, 2022). A busca bibliográfica exploratória sobre animais desaparecidos e soluções digitais foi realizada no portal SBC OpenLib (SOL) e nas ferramentas Crossref e OpenAlex, com termos como ``animais desaparecidos Brasil'', ``animais perdidos aplicativo'' e ``FugaPet-Rotas''. A busca não foi sistemática nem pretendeu esgotar a produção nacional; por isso, a comparação com trabalhos relacionados, como o de Souza et al. (2022), discutido no Capítulo 2, limita-se às fontes consultadas. A análise comparativa, de caráter intencional e não exaustivo, considerou três plataformas voltadas a animais perdidos e encontrados e o Facebook, incluído por seu uso pelos participantes para divulgação dos casos.

Em relação à participação e à proteção dos dados, o texto inicial do formulário apresentou os objetivos acadêmicos da pesquisa e informou seu caráter voluntário; o envio das respostas foi considerado manifestação de consentimento. O formulário não solicitou nomes e estava configurado para não coletar endereços de e-mail. As respostas foram analisadas de forma agregada, buscando preservar a privacidade dos participantes, em consonância com os princípios da Lei Geral de Proteção de Dados Pessoais (Brasil, 2018). Como os questionários consistiram em consultas de opinião a participantes não identificados, o estudo enquadra-se na hipótese de dispensa de registro e avaliação pelo sistema CEP/Conep prevista no art. 1º, parágrafo único, inciso I, da Resolução CNS nº 510/2016 (Brasil, 2016). Os mesmos cuidados foram adotados na avaliação de usabilidade descrita no Capítulo 8.

\subsection{Pesquisa de Campo}

A coleta foi realizada entre 1º e 11 de setembro de 2026 por meio de um questionário digital hospedado no Google Forms, divulgado em um grupo do Facebook relacionado à proteção animal. A participação era aberta a pessoas interessadas no tema, sem exigência de posse ou experiência prévia com animais. Vinte pessoas responderam ao formulário, composto por 18 perguntas. A questão sobre os canais de divulgação foi incluída após o início da coleta e recebeu 19 respostas válidas. As perguntas do instrumento e os gráficos complementares estão no Apêndice B (Seção B.1). A amostra foi não probabilística, por conveniência, e formada por pessoas que optaram voluntariamente por responder ao questionário. Por isso, os resultados descrevem apenas o grupo consultado e não permitem generalização estatística para a população brasileira; além disso, o recrutamento em um grupo ligado à proteção animal pode ter atraído pessoas com maior interesse no tema.

O levantamento investigou práticas de identificação de animais e as experiências e dificuldades dos participantes em situações de desaparecimento. As respostas sobre o uso de microchips estão apresentadas na Figura~\ref{fig:adocao_microchip}.

\begin{figure}[H]
    \centering
    \caption{Adesão ao uso de microchip de identificação em animais domésticos}
    \label{fig:adocao_microchip}
    \includegraphics[width=0.52\textwidth]{imagens/01_pesquisa/grafico_adocao_microchip.png}
    \fonte{Dados da pesquisa (Google Forms, 2026)}
\end{figure}

Na Figura~\ref{fig:adocao_microchip}, 70\% dos participantes (14 respostas) declararam que seus animais não possuem microchip de identificação eletrônica, 10\% (2 respostas) não souberam informar e 20\% (4 participantes) relataram possuir o dispositivo.

Em Pelotas, a Lei municipal nº 7.494/2025 estabelece a microchipagem obrigatória apenas em situações específicas (Pelotas, 2025). O microchip é um identificador que precisa ser lido por equipamento apropriado e não funciona como rastreador GPS.


O questionário também investigou o uso rotineiro de coleiras de identificação e dispositivos de rastreamento, conforme apresentado na Figura~\ref{fig:uso_coleira}.

\begin{figure}[H]
    \centering
    \caption{Utilização de coleiras de identificação em animais domésticos}
    \label{fig:uso_coleira}
    \includegraphics[width=0.68\textwidth]{imagens/01_pesquisa/grafico_uso_coleira.png}
    \fonte{Dados da pesquisa (Google Forms, 2026)}
\end{figure}

Na Figura~\ref{fig:uso_coleira}, 75\% dos participantes (15 respostas) indicaram que seus animais domésticos não utilizam coleira no dia a dia; 15\% (3 respostas) relataram o uso de plaquinha de identificação visual e 10\% (2 respostas), de coleira com dispositivo de rastreamento.


Em seguida, investigou-se o comportamento dos participantes em episódios de desaparecimento, identificando os canais usados no momento do sumiço, conforme demonstrado na Figura~\ref{fig:pesquisa_meios}.

\begin{figure}[H]
    \centering
    \caption{Meios de divulgação para animais perdidos}
    \label{fig:pesquisa_meios}
    \includegraphics[width=0.68\textwidth]{imagens/01_pesquisa/figura01_canais_divulgacao.png}
    \fonte{Dados da pesquisa (Google Forms, 2026)}
\end{figure}

Os resultados da Figura~\ref{fig:pesquisa_meios}, na qual era permitida a seleção de múltiplos canais simultâneos, foram calculados sobre 19 respostas válidas. Nesse grupo, 100\% (19 respostas) relataram recorrer a redes sociais generalistas, como Instagram e Facebook, para divulgar animais perdidos. Em seguida, apareceram os grupos de troca de mensagens no WhatsApp, com 73,7\% (14 respostas), a afixação de cartazes impressos em postes e comércios, com 57,9\% (11 respostas), o contato presencial com vizinhos, com 42,1\% (8 respostas), e a comunicação direta com clínicas veterinárias, com 26,3\% (5 respostas).


Para compreender a vivência prática da comunidade frente à problemática dos desaparecimentos, os participantes foram questionados sobre experiências prévias de perda ou localização de animais desorientados nas ruas, conforme exibido na Figura~\ref{fig:experiencia_perda}.

\begin{figure}[H]
    \centering
    \caption{Experiência prévia com perda ou resgate de animais}
    \label{fig:experiencia_perda}
    \includegraphics[width=0.68\textwidth]{imagens/01_pesquisa/figura02_experiencia_perda.png}
    \fonte{Dados da pesquisa (Google Forms, 2026)}
\end{figure}

A Figura~\ref{fig:experiencia_perda} mostra que 18 dos 20 participantes (90\%) já haviam vivenciado a perda de um animal de estimação. Considerando o total da amostra, 13 participantes (65\%) relataram ter reencontrado o animal e 5 (25\%) não o encontraram. Um participante (5\%) informou ter encontrado ou resgatado um animal desorientado e procurado seu tutor, enquanto outro (5\%) afirmou já ter ajudado amigos ou conhecidos em buscas, sem ter vivenciado pessoalmente a perda.


Na questão sobre dificuldades nos meios atuais, respondida por 19 pessoas, todas apontaram a perda de alcance das publicações após algumas horas; 17 (89,5\%) indicaram a ausência de mapas interativos e a insuficiência de informações detalhadas, e 14 (73,7\%) relataram preocupação com a exposição de telefones pessoais.

As questões sobre funcionalidades tiveram 20 respostas. Todos avaliaram positivamente a utilidade do mapa (notas 4 e 5), e 18 (90\%) atribuíram nota máxima à relevância dos alertas no celular. Sobre o registro de relatos rápidos no mapa, 18 (90\%) responderam ``sim, com certeza'' e 2 (10\%), ``provavelmente sim''. A geração automática de cartazes recebeu avaliação positiva de todos, e o chat privativo, de 19 participantes (95\%). A centralização de anúncios foi apontada como principal diferencial por 9 participantes (45\%). Os resultados completos estão nos gráficos das Figuras~B2 a B14 do Apêndice~B. Esses resultados orientaram a definição dos requisitos apresentados no Capítulo 4.

\subsection{Análise Comparativa de Sistemas Similares}

Para levantar funcionalidades e limitações relevantes à definição dos requisitos, foram examinadas as interfaces e os recursos disponíveis em quatro plataformas: \textit{Viu Meu Pet}, \textit{PawBoost}, \textit{AlertPet} e \textit{Facebook}. A seleção foi intencional e não pretende representar todas as soluções existentes. O Facebook foi incluído por ter sido citado pelos participantes como canal de divulgação. A comparação considerou recursos relacionados à publicação, busca, localização, comunicação e identificação. No Quadro~1, as marcações indicam a presença de recursos próprios e estruturados, observados nas versões consultadas; possibilidades de obter resultado semelhante por busca genérica, notificações comuns ou ações manuais não foram contabilizadas como funcionalidades específicas. As informações sobre recursos pagos foram registradas de forma descritiva, sem constituir um critério de comparação sistemática.


\subsubsection{Viu Meu Pet}

O primeiro sistema analisado foi o \textit{Viu Meu Pet}. A Figura~\ref{fig:viumeupet} apresenta a página inicial da plataforma na versão consultada, com os botões de acesso aos principais recursos.

\begin{figure}[H]
    \centering
    \caption{Interface da plataforma Viu Meu Pet}
    \label{fig:viumeupet}
    \capturaTela{imagens/01_pesquisa/figura03_viumeupet.png}
    \fonte{Viu Meu Pet (2026)}
\end{figure}

\noindent \textit{Descrição geral e interface.} O \textit{Viu Meu Pet} é um serviço brasileiro voltado a animais perdidos e encontrados, disponível em navegador e aplicativo móvel (Figura~\ref{fig:viumeupet}). Na tela consultada, as ações para informar a perda ou o encontro de um animal aparecem em botões destacados, acompanhados de métricas públicas de reencontros e avisos institucionais.

\vspace{0.25cm}
\noindent \textit{Recursos observados.} Na versão consultada, a plataforma oferece gerador de cartazes com código QR para impressão ou compartilhamento digital, mas não foi identificada a sugestão de correspondência entre anúncios.

\vspace{0.25cm}
\noindent \textit{Aspectos observados na análise.} Na análise realizada, os anúncios estavam organizados em listagens por município e estado, sem mapa interativo que apresentasse a localização dos animais anunciados. Também não foram identificados recursos para registrar relatos georreferenciados vinculados a uma publicação. Algumas opções de destaque e notificação dependem de planos pagos.


\subsubsection{PawBoost}

Em seguida, examinou-se o \textit{PawBoost}, plataforma internacional cuja página inicial é apresentada na Figura~\ref{fig:pawboost}.

\begin{figure}[H]
    \centering
    \caption{Interface da plataforma PawBoost}
    \label{fig:pawboost}
    \capturaTela{imagens/01_pesquisa/figura04_pawboost.png}
    \fonte{PawBoost (2026)}
\end{figure}

\noindent \textit{Descrição geral e interface.} O \textit{PawBoost} é uma plataforma internacional voltada à divulgação de animais perdidos e encontrados, disponível em portal web e aplicativos para Android e iOS (Figura~\ref{fig:pawboost}). Na tela consultada, a navegação apresenta atalhos para as principais áreas, modelos de anúncios para compartilhamento e um formulário para iniciar uma publicação.

\vspace{0.25cm}
\noindent \textit{Recursos observados.} Entre os recursos apresentados pela plataforma estão o compartilhamento de alertas em redes sociais, o envio de notificações por e-mail a usuários inscritos em uma região e a geração de anúncios padronizados.

\vspace{0.25cm}
\noindent \textit{Aspectos observados na análise.} Para o contexto brasileiro, a organização territorial baseada em códigos postais dos Estados Unidos (\textit{ZIP Codes}) pode dificultar a localização dos casos. Na versão consultada, não foram identificados mapa público para localizar os animais anunciados, registro georreferenciado de relatos ou recurso de identificação por QR code para coleiras. A plataforma também oferece opções pagas de divulgação.


\subsubsection{AlertPet}

O terceiro sistema analisado foi o \textit{AlertPet}. A Figura~\ref{fig:alertpet} mostra a página da plataforma voltada à contratação de campanhas de divulgação.

\begin{figure}[H]
    \centering
    \caption{Interface da plataforma AlertPet}
    \label{fig:alertpet}
    \capturaTela{imagens/01_pesquisa/figura05_alertpet.png}
    \fonte{AlertPet (2026)}
\end{figure}

\noindent \textit{Descrição geral e interface.} O \textit{AlertPet} é uma plataforma brasileira acessada pelo navegador, que oferece campanhas patrocinadas de divulgação geográfica (Figura~\ref{fig:alertpet}). Na página consultada, aparecem informações sobre a contratação de campanhas, exemplos de anúncios e imagens promocionais com mapas.

\vspace{0.25cm}
\noindent \textit{Recursos observados.} A plataforma oferece a criação de peças visuais para campanhas e a contratação de divulgação geolocalizada em redes sociais, com relatórios de visualizações e engajamento.

\vspace{0.25cm}
\noindent \textit{Aspectos observados na análise.} Na versão consultada, a divulgação depende da contratação de campanhas. Não foram identificados aplicativo móvel próprio, mapa público para localizar os animais anunciados, registro colaborativo georreferenciado, canal privado de mensagens ou identificação por QR code para coleiras.


\subsubsection{Facebook}

Por fim, considerou-se o \textit{Facebook}, citado pelos participantes como canal de divulgação. A Figura~\ref{fig:facebook} apresenta a página de um grupo público dedicado à proteção animal.

\begin{figure}[H]
    \centering
    \caption{Página do grupo SOS Animais no Facebook}
    \label{fig:facebook}
    \capturaTela{imagens/01_pesquisa/figura08_facebook_grupo_sos_animais.png}
    \fonte{Captura de tela do grupo público SOS Animais no Facebook, realizada pelo autor em 2026.}
\end{figure}

\noindent \textit{Descrição geral e interface.} O \textit{Facebook} é uma rede social generalista utilizada pelos participantes para divulgar casos de animais perdidos em grupos comunitários, páginas de proteção animal e perfis pessoais. A Figura~\ref{fig:facebook} apresenta a página do grupo público SOS Animais, com informações do grupo, área para publicação, feed de discussões com textos e fotografias e painel de mídia recente. As publicações permitem interação entre os membros, mas não são estruturadas como registros com campos próprios para localização e avistamentos.

\vspace{0.25cm}
\noindent \textit{Recursos observados.} Na questão sobre canais de divulgação, os 19 participantes que responderam indicaram o uso de redes sociais. Grupos e publicações permitem compartilhar textos e imagens e receber interações de outros usuários sem custo para a publicação orgânica.

\vspace{0.25cm}
\noindent \textit{Aspectos observados na análise.} Como rede social generalista, o Facebook não organiza as publicações sobre animais em uma base estruturada com filtros geográficos ou recursos para localizar os animais anunciados em um mapa. Os relatos acrescentados nos comentários não ficam vinculados a campos próprios de local e horário. Na questão sobre alcance, os 19 participantes relataram que as publicações perdem visibilidade após algumas horas. Além disso, 14 dos 19 respondentes manifestaram preocupação com a exposição de telefones pessoais. Essa preocupação não significa que a divulgação do telefone seja obrigatória, pois a plataforma também oferece mensagens privadas.


\begin{table}[H]
\centering
\phantomsection\hypertarget{quadro:similares}{}\label{quadro:similares}
\footnotesize
\renewcommand{\arraystretch}{1.2}
\setlength{\tabcolsep}{3.5pt}
\begin{tabular}{|>{\raggedright\arraybackslash}p{6.0cm}|>{\centering\arraybackslash}p{2.0cm}|>{\centering\arraybackslash}p{2.0cm}|>{\centering\arraybackslash}p{2.0cm}|>{\centering\arraybackslash}p{2.0cm}|}
\hline
\multicolumn{5}{|c|}{\textbf{Quadro 1: Análise Comparativa de Funcionalidades entre Sistemas Similares}} \\ \hline
\textbf{Funcionalidade} & \textbf{Viu Meu Pet} & \textbf{PawBoost} & \textbf{AlertPet} & \textbf{Facebook} \\ \hline
Cadastrar animal com fotos e descrição & \certo & \certo & \certo & \certo \\ \hline
Disponibilizar aplicativo móvel para busca de animais & \certo & \certo & \errado & \errado \\ \hline
Gerar cartazes digitais de busca & \certo & \certo & \certo & \errado \\ \hline
Filtrar animais por espécie e município & \certo & \certo & \certo & \errado \\ \hline
Enviar alertas sobre publicações relevantes & \certo & \certo & \certo & \errado \\ \hline
\textbf{Visualizar animais em mapa interativo} & \errado & \errado & \errado & \errado \\ \hline
\textbf{Registrar relatos de avistamento com localização} & \errado & \errado & \errado & \errado \\ \hline
Permitir comunicação privada entre usuários & \errado & \errado & \errado & \certo \\ \hline
\textbf{Traçar rota até o local informado na publicação} & \errado & \errado & \errado & \errado \\ \hline
\textbf{Sugerir correspondências entre animais perdidos e encontrados} & \errado & \errado & \errado & \errado \\ \hline
\end{tabular}
\vspace{-0.2cm}
\fonte{Autoria própria (2026)}
\end{table}


A síntese do Quadro~1 e as respostas obtidas na pesquisa indicam que, entre as plataformas examinadas, não foi identificada uma solução que reúna todos os recursos considerados neste trabalho. A partir dessa análise, destacam-se características da proposta do WeFIND:

O primeiro diferencial é a organização geográfica das publicações: enquanto as versões consultadas do Viu Meu Pet e do Facebook apresentam anúncios em listas ou linhas do tempo, o WeFIND também permite localizar os animais anunciados em um mapa interativo, aplicar filtros por raio de distância e traçar uma rota até o local associado à publicação.

O segundo aspecto é a proposta de acesso às funções principais sem cobrança por anúncio, incluindo filtros e sugestões automáticas de correspondência entre animais perdidos e encontrados. Essa correspondência não foi identificada nas plataformas analisadas. O modelo também contrasta com campanhas pagas de divulgação disponíveis em algumas delas.

Outro aspecto é o canal de mensagens associado às publicações, que permite conversar pelo próprio WeFIND sem incluir o telefone no anúncio público. Redes sociais como o Facebook também oferecem mensagens privadas, mas esse contato ocorre fora de um fluxo estruturado de acompanhamento da publicação.

Por fim, o registro de avistamentos vinculado a uma publicação permite organizar relatos com local, horário e foto durante a busca, sem pressupor o trajeto contínuo do animal.

A articulação entre as necessidades identificadas, os requisitos de software e as funcionalidades implementadas encontra-se estruturada na Matriz de Rastreabilidade de Requisitos e Verificações, apresentada no Apêndice~D. A matriz também registra as verificações funcionais realizadas pelo autor; a avaliação de usabilidade com participantes é descrita separadamente no Capítulo~8.

\subsection{Processo de Desenvolvimento de Software}

Como o aplicativo foi desenvolvido em entregas funcionais sucessivas, verificadas em dispositivos móveis e ajustadas antes do avanço para incrementos seguintes, adotou-se um ciclo de vida iterativo e incremental (Valente, 2023).

O processo também se apoiou em princípios do desenvolvimento ágil de software, como a entrega de software funcional e a adaptação a mudanças, articulados ao ciclo iterativo e incremental.

Para estruturar essas premissas no ciclo prático de trabalho, combinou-se o modelo de ciclo de vida iterativo e incremental com o método visual de gerenciamento de tarefas Kanban, operacionalizado em plataforma em nuvem para articular a evolução das funcionalidades desde a concepção dos requisitos até a validação por testes.


\subsubsection{Abordagem Ágil e Ciclo de Vida Incremental}

Para o desenvolvimento da aplicação, adotou-se o modelo de ciclo de vida incremental e iterativo (Valente, 2023). Sob essa abordagem, o sistema não é produzido de uma só vez em um bloco único; pelo contrário, o projeto foi dividido em pequenas entregas chamadas de incrementos funcionais, cada um planejado, desenvolvido e testado no dispositivo móvel antes do início da fase seguinte.

A evolução do aplicativo organizou-se em cinco incrementos funcionais:

A Figura~\ref{fig:ciclo_incremental_wefind} ilustra o modelo incremental de forma geral: cada entrega percorre etapas de desenvolvimento e acrescenta funcionalidades ao produto. No WeFIND, a evolução ocorreu por meio dos cinco incrementos funcionais descritos a seguir.

\begin{figure}[H]
    \centering
    \caption{Representação do modelo incremental de desenvolvimento de software}
    \label{fig:ciclo_incremental_wefind}
    \includegraphics[width=0.82\textwidth]{imagens/03_arquitetura_design/figura13_ciclo_iterativo_incremental.png}
    \fonte{Kolb (2013), que atribui a figura a Pressman (2010)}
\end{figure}

O primeiro incremento dedicou-se ao acesso e à gestão de usuários, compreendendo a criação de conta, a autenticação e a visualização do perfil pessoal.

O segundo incremento contemplou o núcleo espacial de busca, envolvendo o mapa interativo, a exibição de animais perdidos e encontrados por proximidade física e a filtragem rápida por espécie.

No terceiro incremento, desenvolveu-se a gestão de publicações e fotos, permitindo cadastrar anúncios de animais perdidos ou encontrados, acompanhados de características físicas e imagens.

O quarto incremento concentrou-se na comunicação e colaboração, integrando a linha do tempo de relatos de avistamentos pela comunidade e o canal de mensagens privativo em tempo real, permitindo o contato entre os participantes sem a necessidade de expor números de telefone no anúncio.

Por fim, o quinto incremento concentrou-se na geração de cartazes padronizados de busca com código QR para compartilhamento ou impressão.

Essa organização favoreceu a identificação antecipada de inconsistências por meio de verificações incrementais em aparelhos celulares. A avaliação sistemática da cobertura dos requisitos e os testes das funcionalidades são detalhados no Capítulo 8.


\subsubsection{Gestão de Tarefas com Kanban}

Para apoiar a organização e o acompanhamento do fluxo de trabalho ao longo dos incrementos, adotou-se o método de gestão visual Kanban. O método fundamenta-se na visualização do fluxo de trabalho e na limitação das tarefas em andamento, o que ajuda a identificar gargalos e a manter um ritmo de entrega compatível com a capacidade de quem executa as atividades (Valente, 2023). No projeto, sua organização foi realizada na plataforma Trello (Trello, 2026).

A rotina de gestão apoiou-se na visualização das atividades e no acompanhamento de seu andamento. No quadro utilizado, as tarefas foram distribuídas em listas que identificam atividades de atualização, testes e validação, desenvolvimento, itens a fazer, ideias futuras e entregas concluídas. Os cartões permitiram registrar e acompanhar as atividades do projeto ao longo de sua execução.

O quadro também reuniu requisitos funcionais e não funcionais, além de oportunidades de aprimoramento identificadas na pesquisa de campo e na análise comparativa. Para a apresentação no Capítulo 4, os requisitos foram consolidados em 18 requisitos funcionais centrais e quatro requisitos não funcionais verificáveis. A Figura~\ref{fig:quadro_kanban_trello} registra a organização do quadro e exemplifica como as listas distinguem atividades em diferentes situações e propostas reservadas para trabalhos futuros.

\begin{figure}[H]
    \centering
    \caption{Organização das atividades do projeto no quadro Trello}
    \label{fig:quadro_kanban_trello}
    \capturaTela{imagens/03_metodologia/quadro_kanban_trello_wefind.png}
    \fonte{Captura de tela do quadro Trello do projeto, realizada pelo autor (2026).}
\end{figure}


\subsubsection{Etapas do Processo de Desenvolvimento}

Para descrever o processo em linhas gerais, o projeto foi organizado em quatro etapas. Cabe registrar que a implementação do aplicativo já estava em andamento quando a pesquisa de campo foi aplicada, em setembro de 2026; por isso, as respostas obtidas foram utilizadas para validar, priorizar e ajustar os requisitos e funcionalidades em desenvolvimento, e não como único ponto de partida do projeto. Essas etapas sintetizam as atividades metodológicas gerais, enquanto os cinco incrementos representam as entregas funcionais do aplicativo. As atividades de requisitos, modelagem e implementação foram retomadas e ajustadas ao longo dos incrementos; a avaliação formal de usabilidade com participantes ocorreu ao final. A primeira etapa correspondeu à concepção e engenharia de requisitos, na qual as respostas obtidas no formulário de campo com a comunidade e as lições aprendidas na análise comparativa de plataformas similares foram organizadas em requisitos e posteriormente consolidadas no escopo central apresentado no Capítulo 4. Como as etapas se sobrepuseram ao longo dos incrementos, essa consolidação ocorreu de forma iterativa.

A segunda etapa envolveu a modelagem e o projeto de interface, compreendendo a construção dos diagramas conceituais da UML no software Astah UML (Astah, 2026), o planejamento estrutural do banco de dados relacional e a elaboração iterativa do protótipo visual e de interação no Figma (Figma, 2026), apresentados nos Capítulos 4, 6 e 7.

A terceira etapa consistiu na implementação e codificação modular do aplicativo móvel, em que as telas, regras de negócio e serviços de sincronização em nuvem foram programados e integrados com apoio do controle de versões no Git e GitHub (Capítulo 6).

Por fim, a quarta etapa compreendeu a avaliação de usabilidade do aplicativo com participantes voluntários. Os participantes foram convidados a realizar cinco tarefas no WeFIND e, em seguida, responderam a um formulário com questões sobre sua experiência, incluindo dez afirmações estruturadas com base na SUS e perguntas abertas. Os procedimentos, os resultados e as limitações dessa avaliação são apresentados no Capítulo 8.

A integração entre a pesquisa de campo, a análise de similares e as práticas de engenharia de software fundamentou a concepção e a implementação da aplicação móvel. O Capítulo 4 apresenta o levantamento detalhado de requisitos e os modelos conceituais do sistema.
